% 2009 年重庆卷文科第 20 题：双曲线右支、圆 x^2+(y-√5)^2=1 与点 A、B、M。
% 来源 PDF：普通高考/2009/2009重庆文.pdf，SHA-256 c4d433586be8144d34c8ee2d4074c6308f53ffe84aa3a8a7ea28bb6f0cf86a39。
% 原页：第 2 页（A3 三栏）第 2 栏上部，书页码 506。
% 重绘依据：按题设双曲线 x^2-y^2/4=1 与圆 x^2+(y-√5)^2=1 等比绘制；
%   A(-√5,0) 为左焦点，M 取右支上使 |MA|+|MB| 最小的点，B 取圆上相应点。
\begin{tikzpicture}[semithick, x=1.2cm, y=1.2cm, every node/.style={inner sep=1.5pt, font=\small}]
  % 坐标轴
  \draw[-stealth] (-3.0,0) -- (3.0,0) node[below] {$x$};
  \draw[-stealth] (0,-0.5) -- (0,4.0) node[left] {$y$};
  \node[below left] at (0,0) {$O$};

  % 双曲线右支
  \draw[domain=-2.3:2.3, samples=120, thick] plot ({sqrt(1+(\x*\x)/4)},{\x});

  % 圆
  \draw[thick] (0,2.236) circle (1);

  % 点与线段
  \coordinate (A) at (-2.236,0);
  \coordinate (B) at (0.707,1.529);
  \coordinate (M) at (1.140,1.096);
  \fill (A) circle (1.4pt);
  \fill (B) circle (1.4pt);
  \fill (M) circle (1.4pt);
  \draw (A) -- (M);
  \draw[dashed] (M) -- (B);

  % 标注
  \node[below left] at (A) {$A$};
  \node[above left] at (B) {$B$};
  \node[right] at (M) {$M$};
\end{tikzpicture}
